\section{Code Quality Score: Worked Examples}
\label{sec:appendix_c_quality}

This appendix provides concrete worked examples demonstrating how the weighted Code Quality metric (Equation~\ref{eq:code_quality}) is evaluated for different Cirq code implementations.

\subsection{Example 1: Optimal Entangled State Code (Score = 1.00)}
The following code snippet successfully prepares a 2-qubit Bell state ($|\Phi^+\rangle$), measures both qubits, and runs a simulation:

\begin{lstlisting}[language=Python]
import cirq

# Initialize qubits
q0, q1 = cirq.LineQubit.range(2)

# Create circuit
circuit = cirq.Circuit(
    cirq.H(q0),
    cirq.CNOT(q0, q1),
    cirq.measure(q0, q1, key='result')
)

# Run simulation
sim = cirq.Simulator()
result = sim.run(circuit, repetitions=1024)
print(result)
\end{lstlisting}

\noindent
\textbf{Sub-metric Evaluation:}
\begin{itemize}
    \item \textbf{SyntaxValid(c) = 1}: The code successfully parses into a valid Python Abstract Syntax Tree (AST) using \texttt{ast.parse()}.
    \item \textbf{CompilationSuccess(c) = 1}: The code compiles and executes successfully without throwing any exceptions.
    \item \textbf{HasMeasurement(c) = 1}: The code applies a measurement gate (\texttt{cirq.measure}) to the qubits, and the resulting circuit contains at least one \texttt{cirq.MeasurementGate}.
    \item \textbf{CircuitNonEmpty(c) = 1}: The circuit contains three operations (H, CNOT, and measurement), which is greater than zero.
    \item \textbf{CorrectImports(c) = 1}: The code contains a valid import statement for \texttt{cirq} and does not attempt to import unsupported or deprecated libraries.
\end{itemize}

\noindent
\textbf{Code Quality Score Calculation:}
\begin{align*}
\text{CodeQuality}(c) &= 0.30(1) + 0.30(1) + 0.20(1) + 0.10(1) + 0.10(1) \\
&= 0.30 + 0.30 + 0.20 + 0.10 + 0.10 \\
&= 1.00
\end{align

\subsection{Example 2: Code with Syntax Errors (Score = 0.10)}
The following code snippet has a syntax error (missing parentheses for the \texttt{print} function in Python 3):

\begin{lstlisting}[language=Python]
import cirq

q = cirq.LineQubit(0)
circuit = cirq.Circuit(
    cirq.H(q),
    cirq.measure(q, key='r')
)
print circuit
\end{lstlisting}

\noindent
\textbf{Sub-metric Evaluation:}
\begin{itemize}
    \item \textbf{SyntaxValid(c) = 0}: The code fails to parse into a Python AST due to the invalid \texttt{print circuit} syntax, raising a \texttt{SyntaxError}.
    \item \textbf{CompilationSuccess(c) = 0}: Since the code has a syntax error, it cannot compile or run.
    \item \textbf{HasMeasurement(c) = 0}: Although the text contains the word \texttt{measure}, AST analysis and compiler execution cannot extract a valid circuit to verify the measurement gate.
    \item \textbf{CircuitNonEmpty(c) = 0}: The circuit object is not initialized or verified.
    \item \textbf{CorrectImports(c) = 1}: The import statement \texttt{import cirq} is present and valid.
\end{itemize}

\noindent
\textbf{Code Quality Score Calculation:}
\begin{align*}
\text{CodeQuality}(c) &= 0.30(0) + 0.30(0) + 0.20(0) + 0.10(0) + 0.10(1) \\
&= 0.10
\end{align*}

\subsection{Example 3: Valid Syntax but No Circuit Operations (Score = 0.70)}
The following code block is syntactically valid and executes without throwing errors, but it does not define or populate any operations on the circuit:

\begin{lstlisting}[language=Python]
import cirq

# Setup simulator but do not define any gates
q = cirq.LineQubit(0)
circuit = cirq.Circuit()
sim = cirq.Simulator()
\end{lstlisting}

\noindent
\textbf{Sub-metric Evaluation:}
\begin{itemize}
    \item \textbf{SyntaxValid(c) = 1}: The code is syntactically correct and parses into an AST.
    \item \textbf{CompilationSuccess(c) = 1}: The code executes successfully without runtime errors.
    \item \textbf{HasMeasurement(c) = 0}: The code does not define any measurement operations.
    \item \textbf{CircuitNonEmpty(c) = 0}: The circuit is empty (contains zero operations).
    \item \textbf{CorrectImports(c) = 1}: The import statement \texttt{import cirq} is present and valid.
\end{itemize}

\noindent
\textbf{Code Quality Score Calculation:}
\begin{align*}
\text{CodeQuality}(c) &= 0.30(1) + 0.30(1) + 0.20(0) + 0.10(0) + 0.10(1) \\
&= 0.30 + 0.30 + 0.00 + 0.00 + 0.10 \\
&= 0.70
\end{align*}
